\chapter{System Implementation}
\label{ch:system-implementation}

% ─────────────────────────────────────────────────────────────────────────────
\section{Scope and approach}

\subsection{Scope of this chapter}
This chapter documents the \emph{system} built around the forecasting models.
That system comprises the orchestration layer, infrastructure choices,
observability stack, and CI/CD pipeline that together let the platform run
reliably in production. It
does not cover the modelling research itself, such as LEAR hyperparameters, DNN
architecture, the asinh price transformation, or the Diebold--Mariano and
Giacomini--White statistical tests (Chapter~\ref{ch:methods-models}). The
exception is where a modelling decision directly motivates or constrains an
MLOps practice. Operational how-to material such as running backfills, adding
new bidding areas, or re-deploying Prefect flows is kept in the repository's
\texttt{RUNBOOK.md} and \texttt{CONTRIBUTING.md} and is not repeated here.

\subsection{Structuring the chapter around Zarour's practices}
The chapter is organised around the nine MLOps practices identified in the
systematic literature review by Zarour et al.\ \cite{zarour2025}. These are
automating CI/CD, orchestrating workflows, ensuring reproducibility, versioning,
collaboration and communication, continuously training and evaluating models,
tracking and logging ML metadata, continuous monitoring, and implementing
feedback loops. Each practice section follows a fixed shape. It gives Zarour's
own framing of the practice and what it requires of a production ML system, the
concrete implementation found in algo-repo, the technology chosen together with
the alternatives that were considered and rejected, and the design rationale
given the single-host MVP scope of the platform. This
structure makes it straightforward to assess which practices are fully
implemented, partially implemented, or deliberately deferred, and to trace
every technology choice back to an explicit justification. The same nine-practice
vocabulary structures the maturity assessment in
Chapter~\ref{ch:results-system}.

% ─────────────────────────────────────────────────────────────────────────────
\section{System at a Glance}

Before examining each Zarour practice in detail, this section presents the
platform as a whole, covering its host-level topology, the end-to-end data plane,
the complete inventory of scheduled deployments, and the technology choices that
underpin every layer.

\subsection{Host-level architecture}

The platform runs entirely within a single Docker Compose stack on one
dedicated host, keeping operational complexity proportionate to the MVP scope.
Six long-lived services form the persistent backbone of that stack. ClickHouse
stores all time-series data and ML artifacts. The Prefect server exposes the
orchestration API and web UI. The Prefect worker polls for scheduled runs and
manages a Docker work pool. Loki receives structured log streams from Python via
\texttt{python-logging-loki}, and Prometheus scrapes service metrics from
ClickHouse and Grafana. Grafana visualises the ClickHouse data, the Loki log
stream, and the Prometheus series through file-provisioned dashboards. When the Prefect worker
dispatches a scheduled flow run, it spawns an ephemeral container from the
\texttt{algo-repo:latest} image. That container reads from and writes to
ClickHouse, calls the ENTSO-E Transparency Platform API, and emits structured
logs back to Loki before being removed automatically on completion.

\begin{figure}[h]
  \centering
  \includegraphics[width=\linewidth]{f1-architecture}
  \caption{Host-level architecture. A single Docker Compose host runs the six long-lived services, namely ClickHouse, the Prefect server and worker, Prometheus, Grafana, and Loki. Flow runs spawn ephemeral \texttt{algo-repo:latest} containers that read and write ClickHouse, fetch from ENTSO-E, and push logs to Loki; Prometheus scrapes service metrics and Grafana queries Prometheus and Loki.}
  \label{fig:f1-architecture}
\end{figure}

\subsection{Data plane}

Data enters the platform through the ENTSO-E Transparency Platform API, where
four fetch functions retrieve energy prices, total-load forecasts, generation
forecasts, and wind and solar forecasts for each configured bidding area.
The raw XML responses are parsed by typed mapper functions that produce
\texttt{Observations} and \texttt{Forecast} domain objects, which are then
ingested into ClickHouse tables \texttt{algo.timeseries} and
\texttt{algo.forecasts} via Arrow-batch inserts. Both tables use
\texttt{ReplacingMergeTree} engines, which deduplicate lazily on merge. Read
paths therefore either append the \texttt{FINAL} keyword to their queries or
call \texttt{drop\_duplicates} on the loaded DataFrame to ensure correctness. A
preprocessor loads this data, pivots it to a daily-wide feature matrix through a
configurable sequence of feature transformers, and feeds it to either the
backtest engine for offline evaluation or the live runner for day-ahead
production forecasting. Predictions and trained model artifacts are persisted to
\texttt{algo.experiment\_runs}, \texttt{algo.experiment\_predictions}, and
\texttt{algo.model\_artifacts}, from where the reconciliation flow back-fills
actuals the following day and the drift-detection flow compares the recent price
distribution against a ninety-day reference window, triggering an alert and an
optional automatic retraining run when significant drift is detected. Figure~\ref{fig:f2-data-plane} draws this
end-to-end path together; it is deferred to the end of the chapter
(\S\ref{sec:si-crosscutting}) so that each stage is introduced first.


\subsection{Operational inventory}

Table~\ref{tab:inventory} lists every Prefect deployment currently registered
in the system, together with its schedule, the flow it executes, and the
Zarour practice it primarily realises.

\begin{table}[H]
  \centering
  \scriptsize
  \caption{Operational inventory of Prefect deployments. Cron strings are UTC; the seven backfill deployments and six forecast deployments share one flow definition apiece.}
  \label{tab:inventory}
  \setlength{\tabcolsep}{4pt}
  \begin{tabularx}{\linewidth}{@{}l l l l X l@{}}
    \toprule
    \textbf{Deployment} & \textbf{Area / model} & \textbf{Cron} & \textbf{Flow} & \textbf{Purpose} & \textbf{Zarour} \\
    \midrule
    entsoe-backfill-dk1 & DK1 & 0 * * * * & \texttt{entsoe\_backfill\_flow} & Hourly ENTSO-E ingest & P2 \\
    entsoe-backfill-dk2 & DK2 & 0 * * * * & {\itshape (as above)} & Hourly ENTSO-E ingest & P2 \\
    entsoe-backfill-nl  & NL  & 0 * * * * & {\itshape (as above)} & Hourly ENTSO-E ingest & P2 \\
    entsoe-backfill-be  & BE  & 0 * * * * & {\itshape (as above)} & Hourly ENTSO-E ingest & P2 \\
    entsoe-backfill-de  & DE  & 0 * * * * & {\itshape (as above)} & Hourly ENTSO-E ingest & P2 \\
    entsoe-backfill-fr  & FR  & 0 * * * * & {\itshape (as above)} & Hourly ENTSO-E ingest & P2 \\
    entsoe-backfill-pl  & PL  & 0 * * * * & {\itshape (as above)} & Hourly ENTSO-E ingest & P2 \\
    \midrule
    daily-forecast-dk1\_lear        & DK1 / LEAR        & 0 6 * * * & \texttt{daily\_forecast\_flow} & Daily refit + forecast & P6 \\
    daily-forecast-de\_lear         & DE / LEAR         & 0 6 * * * & {\itshape (as above)} & Daily refit + forecast & P6 \\
    daily-forecast-dk1\_dnn         & DK1 / DNN         & 0 6 * * * & {\itshape (as above)} & Daily refit + forecast & P6 \\
    daily-forecast-de\_dnn          & DE / DNN          & 0 6 * * * & {\itshape (as above)} & Daily refit + forecast & P6 \\
    daily-forecast-dk1\_transformer & DK1 / Transformer & 0 6 * * * & {\itshape (as above)} & Daily refit + forecast & P6 \\
    daily-forecast-de\_transformer  & DE / Transformer  & 0 6 * * * & {\itshape (as above)} & Daily refit + forecast & P6 \\
    \midrule
    reconcile-predictions & all & 0 18 * * * & \texttt{reconcile\_predictions\_flow} & Day-after reconciliation & P7, P8 \\
    \midrule
    drift-detection-dk1 & DK1 & 0 19 * * * & \texttt{drift\_detection\_flow} & Drift: 90d ref vs 7d (Evidently) & P8, P9 \\
    drift-detection-de  & DE  & 0 19 * * * & {\itshape (as above)} & Drift: 90d ref vs 7d (Evidently) & P8, P9 \\
    \bottomrule
  \end{tabularx}
\end{table}

\subsection{Technology stack}

Table~\ref{tab:tech-stack} lists each tool and its role; the rationale for each
choice is given in the practice section where the tool appears.

\begin{table}[H]
  \centering
  \caption{Technology in use and its role. The rationale for each choice is given in the practice section where the tool appears; tools evaluated but not adopted (MLflow, Kubernetes, DVC) are discussed there too.}
  \label{tab:tech-stack}
  \begin{tabular}{@{}l l@{}}
    \toprule
    \textbf{Tool} & \textbf{Role} \\
    \midrule
    Prefect          & Orchestration \\
    ClickHouse       & OLAP store + model-artifact blob \\
    Docker Compose   & Container runtime + infrastructure-as-code \\
    Loki             & Log aggregation \\
    Prometheus       & Metrics time-series store \\
    Grafana          & Dashboards and file-based alerts \\
    Evidently AI     & In-process drift detection \\
    W\&B             & Research + production tracking \\
    GitHub Actions   & CI/CD \\
    Tailscale        & Deploy-host network access \\
    Optuna           & Hyperparameter search (TPE) \\
    \bottomrule
  \end{tabular}
\end{table}

% ─────────────────────────────────────────────────────────────────────────────
\section{Implementing Zarour's Nine MLOps Practices}

The nine sections that follow map every MLOps practice identified by Zarour et
al.\ (Chapter~\ref{ch:background}) to a concrete subsystem in algo-repo, in
Zarour's order from P1 to P9. Each follows the same shape, covering the
practice's framing, the implementation found in this repository, the technology
chosen and its rejected alternatives, and the design rationale scoped to a
single-host MVP.
Where a practice is only partially realised, the limitation is stated
explicitly and cross-referenced to the maturity assessment in
Chapter~\ref{ch:results-system}.

\subsection{Automating CI/CD}

Automating continuous integration and delivery requires that every proposed
change be validated before it reaches \texttt{main}, and that an accepted change
then propagate to the running system without manual intervention
\cite{zarour2025}. In algo-repo this is a two-phase GitHub Actions pipeline
(Figure~\ref{fig:f3-cicd}).

On a pull request, three checks must pass before a merge is allowed.
\texttt{pr\_test} runs the full pytest suite mirroring \texttt{src/},
\texttt{pr\_ruff-lint} runs the \texttt{ruff} linter, and \texttt{pr\_smoke} runs
an end-to-end daily-forecast and reconciliation on a throwaway Compose stack with
synthetic data. A path filter skips the expensive smoke build for
documentation-only changes.

On merge to \texttt{main}, \texttt{cd\_deploy} reaches the production host and,
gated by which paths changed, applies any new ClickHouse migrations, rebuilds the
\texttt{algo-repo:latest} image with the commit SHA baked in, and reruns
\texttt{deploy\_ingestion.py} and \texttt{deploy\_forecast.py} to register the
updated Prefect deployments. It finally triggers \texttt{daily-forecast-dk1\_lear}
and polls until it completes, confirming the deployed image produces a valid
forecast on the live host. A parallel job restarts Grafana when its provisioning
changes.

\begin{figure}[h]
  \centering
  \includegraphics[width=0.78\linewidth]{f3-cicd}
  \caption{The CI/CD pipeline. On a pull request, three checks (\texttt{pr\_test}, \texttt{pr\_ruff-lint}, \texttt{pr\_smoke}) must pass before a merge. On merge to \texttt{main}, \texttt{cd\_deploy} reaches the host over Tailscale, rebuilds the image, redeploys the Prefect flows, and runs a smoke forecast.}
  \label{fig:f3-cicd}
\end{figure}

The deploy host is reached only over a Tailscale mesh. The runner joins the
tailnet for the duration of the workflow with an ephemeral key, so the host
never exposes an SSH port to the public internet. GitHub Actions was chosen
over Jenkins or CircleCI because it is co-located with the repository and needs
no separately managed runner.

\subsection{Orchestrating workflows}

Orchestrating workflows, in Zarour et al.'s framing, is the practice of
expressing the recurring data and model tasks as scheduled, observable,
retryable pipelines rather than as cron-driven scripts, so that failures are
isolated, retried, and surfaced rather than silently lost \cite{zarour2025}.

All recurring work on the platform is expressed as four Prefect flows executed
by a Docker work pool (Table~\ref{tab:inventory}). These are ingestion, daily
forecast, reconciliation, and drift detection. Scheduling is declarative. Each deployment in
Table~\ref{tab:inventory} attaches a cron schedule to a flow, and the worker
spawns an ephemeral \texttt{algo-repo:latest} container per run. The ingestion
flow is parameterised by area, a single flow definition with one deployment per
bidding zone, so adding a market is a one-line deployment change with no
flow-code edit.
Within a flow, granularity is one task per data type, covering prices, load,
generation, and wind/solar. Tasks run sequentially to respect the ENTSO-E rate limit and
each carries \texttt{retries=2} with a 60-second delay, so one data type
failing does not abort the others. The ingestion flow re-ingests a rolling
48-hour lookback window on every run. Because both ClickHouse tables are
\texttt{ReplacingMergeTree}, re-ingesting overlapping data is always safe.

A deliberate design choice keeps the fetch-and-insert logic in a shared ingest
service with no Prefect dependency, so that both the data-backfill CLI and the
ingestion flow import the same code path rather than duplicating it.
Prefect was chosen over Airflow and Dagster. At single-host scale its Python-native
flow definitions and Docker work pool are sufficient, and it avoids the
operational weight of an Airflow scheduler or a Dagster daemon and database.

\subsection{Ensuring reproducibility}

Reproducibility in the context of MLOps requires that a past model result can
be reconstructed from the same code version, data window, hyperparameters, and
\emph{runtime environment}, without relying on implicit host state
\cite{zarour2025}. Zarour et al.\ identify it as one of the most frequently
cited challenges in practitioner surveys, and treat containerisation of the
runtime as one of its principal enablers.

The platform addresses reproducibility along three axes. First, the runtime is
fully containerised and the topology is infrastructure-as-code. The entire
operational stack, comprising all six long-lived services with their volumes,
environment, and network, is declared in \texttt{docker-compose.yml}; the
flow-run and worker containers are built from \texttt{Dockerfile} and
\texttt{Dockerfile.worker}; and Grafana's provisioning, the alert rules, and the
Prometheus scrape config are all file-mounted from the repository. The
infrastructure can therefore be rebuilt from the repository alone, with no
hand-applied host commands. Second, every row in \texttt{algo.experiment\_runs} records its
git SHA and the data window it used, which together fix the code version and the
data behind any run. The \texttt{GIT\_SHA} is baked into the image at
build time rather than resolved at runtime, because the deploy host holds only
the built image and not the source tree, and local runs emit an empty string so
a dirty tree never stamps a misleading SHA. Third, the modelling pipeline is itself deterministic given the
config. Calibration windows and the \texttt{AsinhScaler} target transform are
fixed in the committed configs, and the DNN ensemble averages
\texttt{n\_dnn\_ensemble} replicates with distinct, enumerated random seeds, so
the trained ensemble is fully described by the config and a deterministic seed
sequence.

Docker Compose was chosen over Kubernetes for the container runtime. At
single-host scope, a control plane, etcd, and Helm charts are unjustified
overhead (\S\ref{sec:si-crosscutting}). Baking the
SHA into the image rather than resolving it at container start makes every
\texttt{algo.experiment\_runs} row self-describing under a deployment model that
has no source checkout on the host.

\subsection{Versioning}

Versioning is the foundational practice under which all code, data schemas, and
configuration that affect model behaviour are tracked in source control,
enabling reproducible reconstruction of any past system state
\cite{zarour2025}. Zarour et al.\ treat it as a prerequisite for every other
practice, because without a shared history of code and configuration changes,
claims about reproducibility or rollback cannot be substantiated.

The platform uses Git with GitHub as the remote, and versions every artefact
that can affect a model run. Python experiment configs in \texttt{configs/} are
first-class Git files whose history records every calibration-window or scaler
change. ClickHouse schema changes live in numbered SQL migration files under
\texttt{src/infrastructure\_clickhouse/migrations/} that are auto-discovered in
sorted order by \texttt{run\_migrations.py} and never modified retroactively.
Grafana dashboard definitions and alert rules, stored as JSON and YAML under
\texttt{grafana/}, are file-provisioned, never edited through the UI, and subject
to code review before taking effect.

DVC was evaluated for data versioning and rejected. ENTSO-E data is replayable
in full from the public API, and combining ClickHouse's
\texttt{ReplacingMergeTree} engine with the \texttt{git\_sha} column on
\texttt{algo.experiment\_runs} provides sufficient provenance without
maintaining a separate binary store. Versioning every artefact that affects
model output means a single \texttt{git log} spans the experiment, the
infrastructure, and the alerting policy together, closing the gap between
operational changes and the audit trail expected of a production ML system.

\subsection{Collaboration and communication}

Collaboration practices let contributors work on a shared ML system without
relying on tribal knowledge \cite{zarour2025}. As a single-developer project this
was the least-emphasised of the nine practices, but the core disciplines were
still followed. Every code change went through a GitHub \emph{pull request} rather
than a direct commit to \texttt{main}, so each change passed the CI gate of P1
(Figure~\ref{fig:f3-cicd}) and left a reviewable record of what changed and why.
GitHub \emph{Issues} tracked outstanding work and the status of each feature
branch, giving a single queue of what was in progress, blocked, or done.

Beyond the version-control workflow, project conventions, repository structure,
and operational procedures are kept in versioned Markdown files, namely
\texttt{CLAUDE.md}, \texttt{docs/REPO\_OVERVIEW.md}, \texttt{docs/RUNBOOK.md}, and
\texttt{docs/CONTRIBUTING.md}, so the design rationale and the onboarding path
for a future contributor live in the repository itself rather than in any one
person's head.

\subsection{Continuously train and evaluate}

Continuous training and evaluation requires that the model be retrained
automatically on a regular cadence so that it reflects the most recent data
distribution, and that its live predictions be continuously scored against
realised outcomes \cite{zarour2025}. Zarour et al.\ distinguish the scheduled,
unconditional refit from the reactive drift-triggered path of P9.

The platform's train-and-evaluate cycle runs on a fixed daily cadence
(Figure~\ref{fig:p6-cadence}). The six \texttt{daily-forecast-*} deployments,
covering LEAR, DNN, and Transformer for both DK1 and DE, execute the
daily-forecast flow at 06:00 UTC. Each run invokes the live runner, which
loads the latest ClickHouse data window, refits the ensemble of
calibration-window models from scratch, and immediately produces a day-ahead
forecast. For DNN configurations
it fits \texttt{n\_dnn\_ensemble} seed replicates and averages them, following
Lago et al.\ \cite{lago2021}. The trained models are serialised to
\texttt{algo.model\_artifacts}. The \emph{evaluate} half of the practice is the
reconciliation flow (\S P7--P8). The day after each forecast, realised prices
are joined to the stored predictions and the same MAE, RMSE, and sMAPE used in
backtesting are computed on the live forecast, so offline and online evaluation
use one metric implementation.

\begin{figure}[h]
  \centering
  \includegraphics[width=0.92\linewidth]{f5-daily-cadence}
  \caption{The daily train-and-evaluate cadence (UTC). At 06:00 every configuration refits from scratch and produces a day-ahead forecast; at 18:00 reconciliation joins the published actuals to the stored predictions and scores them; at 19:00 drift detection runs. The cycle repeats every day.}
  \label{fig:p6-cadence}
\end{figure}

A full refit from scratch on each run was chosen over warm-start learning. The
daily data increment is small relative to the calibration windows, so a clean
refit is cheap and avoids accumulated approximation error. Keeping training and
prediction in one flow means a single task-level retry covers both phases and
the served model is always consistent with the forecast it produced.

\subsection{Track and log ML metadata}

Tracking and logging ML metadata is the practice of recording the full
provenance of every model run in a queryable store, so that past results can be
reproduced and compared without ad-hoc file management \cite{zarour2025}. That
provenance covers the hyperparameters, data windows, performance metrics, and the
model artefact itself.

This is handled in ClickHouse. \texttt{algo.experiment\_runs} stores a
self-describing row per run, holding its training duration, data window, git SHA,
and a link to the stored model. \texttt{algo.experiment\_predictions} holds one row
per period per day for reconciliation, \texttt{algo.model\_artifacts} stores the
serialised models as blobs, and \texttt{algo.hyperparam\_trials} persists every
Optuna trial for resumable search. On the research surface, W\&B hosts two
projects, \texttt{algo-research} for backtests and sweeps and \texttt{algo-prod}
for production reconciliation; \texttt{WANDB\_MODE=disabled} no-ops it in CI and
local development.

A filesystem model store was rejected. Storing artefacts as blobs in
\texttt{algo.model\_artifacts} keeps them attached to the experiment row rather
than in a separately managed path. MLflow was evaluated and deferred until the
experiment-config count exceeds ten. At the current scale ClickHouse and W\&B
together cover the audit-log and research-surface roles without adding a service
to the Compose stack. Using one OLAP store for predictions, metrics, and
artefacts makes day-after reconciliation a single query, and the
\texttt{ReplacingMergeTree(created\_utc)} engine makes the reconciliation update
a safe re-insert rather than a mutable \texttt{UPDATE}.

\subsection{Continuous monitoring}

Continuous monitoring requires instrumenting the production system across at
least three axes, namely data quality, model performance, and system health, so
that degradation is detected automatically and surfaced without manual
inspection \cite{zarour2025}.

Structured log events from every Prefect flow and infrastructure module are
pushed directly to Loki via a \texttt{python-logging-loki} handler, with no
Promtail sidecar.
Three Grafana dashboards auto-provision from \texttt{grafana/dashboards/}. The
Data Quality dashboard covers coverage and freshness of \texttt{algo.timeseries}
and \texttt{algo.forecasts}. The Operational dashboard presents the Loki log
stream filtered by \texttt{job=algo-repo}. The System Health dashboard is
backed by Prometheus, which scrapes ClickHouse's native \texttt{/metrics}
endpoint, Grafana's \texttt{/metrics} endpoint, and itself on a 15-second
cadence and retains 120 days of series. The model-quality metrics MAE, RMSE, and
sMAPE, computed per reconciled day, area, and config, do not live in Grafana.
They are owned by the W\&B \texttt{algo-prod} project, which the W\&B tracker
updates from the reconciliation flow.
The split is deliberate. Grafana hosts platform-and-data-health concerns, W\&B
hosts model-quality concerns, and both read from the same ClickHouse tables.

File-based alert rules in \texttt{grafana/provisioning/alerting/rules.yaml}
define four conditions. \texttt{data-staleness} fires when \texttt{algo.timeseries}
has received no new row in twelve hours. \texttt{drift-detected-today} fires when
any row in \texttt{algo.drift\_reports} for the current day carries
\texttt{drift\_detected = 1}. \texttt{reconciliation-missed} fires when no
\texttt{reconcile-predictions} run has completed in the last 36 hours.
\texttt{clickhouse-down} fires when the Prometheus \texttt{up\{job="clickhouse"\}}
liveness metric is zero. In addition, an accuracy-degradation check in the
reconciliation flow raises an alert whenever today's MAE exceeds 1.5\,$\times$ the
30-day median over \texttt{algo.experiment\_runs}.

Loki was chosen over the ELK stack because it is a single binary with
substantially lower memory and operational overhead, proportionate to a
single-host deployment, and Promtail was rejected in favour of direct Python
push, which removes a moving part from the logging pipeline. Prometheus sits
alongside Loki as the metrics tier to Loki's logs tier, with Grafana the single
dashboard and alerting surface above both. The three monitoring axes match
Zarour's P8 decomposition and ensure an alert is raised at whichever layer first
shows a problem.

\subsection{Implementing feedback loops}

The feedback-loop practice requires that the production system detect when its
inputs or outputs have shifted sufficiently to degrade prediction quality and
respond automatically, rather than waiting for a human to notice on a dashboard
\cite{zarour2025}. Zarour et al.\ identify the absence of an automated feedback
loop as a common maturity gap.

Drift detection runs in-process inside the \texttt{drift-detection-dk1} and
\texttt{drift-detection-de} deployments via Evidently's \texttt{DataDriftPreset},
which compares a 90-day reference window against a 7-day current window on each
area's price series and auto-selects the statistical test by sample size, using
a Kolmogorov--Smirnov test for small samples and the Wasserstein distance for
larger ones. Each drift result, comprising its score, threshold, the test
applied, and the boolean \texttt{drift\_detected} flag, is written to
\texttt{algo.drift\_reports}, and a detected event posts a warning to Loki and
raises an alert. When a config carries \texttt{retrain\_on\_drift=True}, the flow
additionally enqueues a triggered-retrain flow, which writes a \texttt{running}
audit row to \texttt{algo.retrain\_events}, calls the daily-forecast flow
synchronously,
and closes the row with a \texttt{completed} or \texttt{failed} status and a
strictly later \texttt{updated\_utc} so the \texttt{ReplacingMergeTree(updated\_utc)}
engine retains the terminal row.

Evidently was deployed in-process rather than as a separate service so the
check inherits the flow's retry policy and is re-run via the Prefect UI without
managing another endpoint. Arize and WhyLabs were rejected as cloud-hosted
dependencies unjustified at single-host scope. Making the retrain trigger
conditional on \texttt{retrain\_on\_drift} keeps every retraining event
auditable. Each triggered run links to the specific \texttt{drift\_report\_id}
that caused it, so the full chain from detected drift to updated model is
recoverable from ClickHouse alone. Whether drift \emph{should} trigger
retraining at all, given the cheap daily refit, is examined in
Chapter~\ref{ch:results-system}.

% ─────────────────────────────────────────────────────────────────────────────
\section{Reproducibility, monitoring, and single-host scope}
\label{sec:si-crosscutting}

With every practice now described, Figure~\ref{fig:f2-data-plane} draws them
together into the single data path the platform implements, from ingestion into
the store, through the forecasting pipeline that reads it, to the daily
reconciliation-and-drift loop that closes back onto the store.

\begin{figure}[H]
  \centering
  \includegraphics[width=\linewidth]{f2-data-plane}
  \caption{Data plane (high level). ENTSO-E data is ingested into ClickHouse; the forecasting pipeline loads it, builds features, and trains and serves the models, writing predictions and artifacts back. The daily reconciliation-and-drift loop reads those predictions, writes back actuals and metrics, and raises alerts or retraining.}
  \label{fig:f2-data-plane}
\end{figure}

\subsection{GIT\_SHA propagation as the reproducibility--monitoring link}

Every row that \texttt{cd\_deploy.yml} produces in \texttt{algo.experiment\_runs}
carries the same \texttt{git\_sha} baked into the image during the deploy-gate
build step, the SHA of the exact commit merged to \texttt{main}. The Grafana
Operational dashboard queries the Loki log stream filtered by
\texttt{job=algo-repo} and presents structured events that include the same SHA,
because every flow container is built from that image. So if the W\&B
\texttt{algo-prod} run for an \texttt{(area, config\_name)} shows an unexpected
MAE spike beginning on a particular day, the Operational dashboard shows the
Loki events from the same period alongside the \texttt{git\_sha} of the image
that produced them, making it possible to \texttt{git log} back to the exact
commit. This links the reproducibility of P3 to the monitoring of P8.
The audit trail for a production anomaly runs from the Grafana panel, through
the \texttt{algo.experiment\_runs} row, to a checkable commit, with no gap.

\subsection{Why a single-host deploy is the right MVP scope}

The entire stack runs on one Docker Compose host, for two reasons. First,
Kubernetes and multi-host orchestration were out of scope for this project. The
operational complexity of a control plane, distributed storage for ClickHouse,
and cross-host log aggregation would have slowed iteration without improving the
platform's ability to deliver daily forecasts. Such an architecture would,
however, genuinely improve robustness, since it removes the single failure domain
and enables rolling deployments and horizontal scaling. It is the natural next
step once a second host is justified. Second, only one compute instance was available,
so a distributed deployment was not an option in any case. On a single host the
blast radius is one failure domain, and rollback is simply rebuilding the image
and restarting the stack.

% Bibliography is included by the master thesis document, not here.
